% Appendix fragment; requires graphicx. Include after \appendix.
% Data: ../evaluations/online_50/{arm_summary,cohort_summary,paired_differences}.csv.
\section{Additional online results}
\label{app:online-details}

Table~\ref{tab:online-results} and Figure~\ref{fig:online-success-time}
expand the main results in Figures~\ref{fig:online-outcomes}
and~\ref{fig:online-model-resources}. The shared geometric implementation and
examples are in Appendix~\ref{app:geometric-abstraction}.

\begin{table*}[t]
  \centering
  \small
  \setlength{\tabcolsep}{5pt}
  \begin{tabular}{llrrrrrr}
    \hline
    Task & Condition & Success & Time (s) & Model (s) & Tokens (k) & Requests & Images \\
    \hline
    Bridge-tower & Base & 12/50 & 198.0 & 55.5 & 112.5 & 7.88 & 7.50 \\
     & Full & 28/50 & 185.3 & 51.8 & 129.5 & 7.32 & 4.28 \\
    \hline
    Ring-on-peg & Base & 49/50 & 115.8 & 90.0 & 214.5 & 11.88 & 10.94 \\
     & Full & 50/50 & 77.2 & 45.6 & 101.3 & 6.90 & 2.72 \\
    \hline
    Tower relocation & Base & 15/50 & 186.4 & 80.3 & 185.9 & 11.58 & 11.34 \\
     & Full & 42/50 & 139.2 & 31.7 & 88.3 & 5.36 & 3.00 \\
    \hline
  \end{tabular}
  \caption{\textbf{Online outcomes and cost.} Success counts completed episodes
  out of 50. Time is mean capped completion time; Model is mean cumulative
  request wall time, including interrupted waits. Tokens are mean recorded
  input plus output per episode (thousands), counting cached input once and
  excluding unknown interrupted usage. Requests counts completed model requests;
  Images counts newly attached images across all request attempts, not local
  camera frames. Both are per-episode means. All means include timeouts; both arms of nine runtime-flagged pairs use
  audited repeats, as detailed below.}
  \label{tab:online-results}
\end{table*}

\begin{figure*}[t]
  \centering
  \includegraphics[width=\textwidth]{figures/success_vs_time.png}
  \caption{\textbf{Success within elapsed time.} Each curve summarizes 50
  selected 200-second episodes, stepping at first verified success. Timeouts
  never become successes at the deadline. Earlier time points are read from the
  same episodes, not separate evaluations at smaller budgets.}
  \label{fig:online-success-time}
\end{figure*}

\begin{figure*}[t]
  \centering
  \includegraphics[width=\textwidth]{figures/paired_effects.png}
  \caption{\textbf{Paired effects and sensitivity.} Full minus Base success
  and capped-time differences for all 50 pairs, the new 45 pairs, and a subset
  excluding whole pairs with historical timing or GPU-query flags. Intervals
  use 10,000 paired-layout percentile bootstrap resamples, seed 1050 per task,
  without multiplicity correction. Positive success and negative time differences
  favor Full. The sensitivity retains 48, 46, and 47 pairs respectively.}
  \label{fig:online-paired-effects}
\end{figure*}

\paragraph{Cohorts and runtime sensitivity.}
The extension was registered after observing the five-pair pilot. On the 45 new
pairs, Base/Full successes are 10/45 versus 24/45 for Bridge-tower, 44/45 versus
45/45 for Ring, and 13/45 versus 39/45 for Relocation; Full's mean capped time is
lower in all three tasks. The model
client changed from 0.155.0-alpha.16.3 to 0.159.2. GPU-runtime and Ring
camera-pacing repairs were recorded as separate execution epochs, with no change
to prompts, skills, geometry fitting, evaluators, or the 200-second budget.

The original runs had four timing violations and sixteen failed dispatch GPU
queries, affecting nine pairs. Both conditions were repeated once on all nine
pairs with unchanged seeds, prompts, skills, evaluators, and budgets; all 18
repeats passed audit. After inspecting the repeats, we amended reporting to use
all 18 in the primary results, including outcomes that became worse. The export
preserves every original result and replacement mapping. The original-attempt
Bridge-tower counts were 12/50 and 27/50; other task totals were unchanged.

Excluding those nine pairs entirely gives Base/Full successes of 12/48 versus
27/48, 45/46 versus 46/46, and 14/47 versus 39/47, with lower Full capped time
on every task (Figure~\ref{fig:online-paired-effects}). Two additional Ring
terminal-image failures were independently verified from saved evidence and
were not rerun. The selected episodes include 90,796 tokens recovered from four
terminal requests using agreeing CLI/session counters; 73 interrupted usages
remain unknown. Superseded attempts used another 2,745,300 recorded tokens and
startup failures 98,036, reported separately from selected-episode means.

\paragraph{Evaluation notes.}
All online experiments use MuJoCo 3.11.0 with a simulated SO-100 robot.
Bridge-tower and Relocation require a five-second stable hold; Ring requires two
seconds. Relocation must preserve or restore both foundations within 5\,mm of
their initial positions. Ring uses the same sensor-driven grasp/lift prefix in
both conditions and returns control before alignment and release. Success and
all costs include this prefix. One selected episode per condition/layout enters the analysis;
paired outcomes can differ, so the primary time means include timeouts rather
than selecting only successful episodes. Model time includes service waiting
and is not isolated inference compute time.
